\section{Related Work}

\paragraph{Segmentation metrics and problem-aware validation.}
Mean intersection over union is widely used for semantic segmentation, but it
compresses class- and pixel-level error structure into one aggregate. Boundary
IoU is designed to increase sensitivity to contour errors, particularly where
standard region overlap is insensitive to object-scale boundary deviations
\cite{cheng2021boundary}. More broadly, validation metrics can produce
misleading comparisons through implementation, aggregation, ranking, or
interpretation choices. Metrics Reloaded recommends selecting complementary
measures from an explicit description of the target problem, and its companion
analysis organizes recurrent metric-related pitfalls
\cite{maierhein2024metrics,reinke2024metrics}. We operationalize one such
concern: different correctly computed criteria can induce incompatible model
preferences. Our finite-panel bound quantifies that incompatibility for the
criteria declared here; it does not identify a universally optimal metric set.

\paragraph{Rank aggregation and social choice.}
Combining several rankings into one is a classical problem. Arrow's
impossibility theorem rules out aggregation rules that satisfy a few natural
axioms simultaneously~\cite{arrow1951social}, and the Kemeny--Young rule selects
the ranking that minimizes total pairwise disagreement with the
voters~\cite{kemeny1959mathematics,young1978consistent}; computing it is NP-hard
in general~\cite{bartholdi1989voting} and has been studied for combining search
engine rankings~\cite{dwork2001rank}. We treat segmentation quality criteria as
voters over models and use the Kemeny objective not to aggregate them, but to
measure how much disagreement any single quality score must incur.

\paragraph{Reliability and uncertainty in segmentation.}
Semantic-segmentation reliability includes distinct tasks such as robustness,
calibration, misclassification detection, and out-of-distribution detection.
A broad architecture study found that architectural progress on robustness
does not imply uniform improvement in uncertainty-based reliability
\cite{dejorge2023reliability}. ValUES advocates controlled validation across
uncertainty applications and emphasizes that score aggregation is itself a
consequential design choice \cite{kahl2024values}. Our scope is narrower: we
compare fixed per-image quality scores for model ranking and failure
discrimination. We do not estimate pixelwise predictive calibration, detect
out-of-distribution inputs, or claim that these tasks are interchangeable.

\paragraph{Safety- and relevance-aware evaluation.}
Driving-segmentation metrics can account for spatial error structure
\cite{cheng2021safety}. Relevance-weighted IoU emphasizes corner-case pixels
\cite{breitenstein2023really}, while Safe mIoU represents class-confusion costs
through a hierarchy \cite{shaik2024iddaw}. These approaches make task
priorities explicit. We similarly declare severity weights and operational
thresholds, but treat them as experimental assumptions rather than validated
harm models. We report policy agreement, operational agreement, severe misses,
and false positives separately so that one endpoint cannot stand in for all
others.

\paragraph{Tail risk, selection, and endpoint choice.}
Conditional value at risk (CVaR) summarizes the worst part of a loss
distribution \cite{rockafellar2000optimization}; minimum component coverage
provides an extreme component-level statistic. Selective-prediction methods
instead evaluate risk as a function of retained coverage and can provide
guarantees under their assumptions \cite{geifman2017selective}. We evaluate
neither abstention nor risk--coverage guarantees. Our question is whether
failure-detection discrimination and model-ranking agreement can be treated as
the same endpoint. The results show why that implication must be tested rather
than assumed: a mixed-failure AUROC gain does not guarantee better policy or
operational ranking.

\paragraph{Architectural diversity.}
The evaluated panels include FCN~\cite{long2015fully}, DeepLabV3
\cite{chen2017rethinking}, MobileNetV3~\cite{howard2019searching},
SegFormer~\cite{xie2021segformer}, UPerNet~\cite{xiao2018unified}, ConvNeXt
\cite{liu2022convnext}, Swin~\cite{liu2021swin}, BEiT~\cite{bao2022beit}, and
Mask2Former~\cite{cheng2022mask2former}. This permits a limited
cross-architecture analysis, but each panel contains only six evaluated
checkpoints. It is not a broad sample of the architecture space.
